% when: xcolor, beamer, tikz, pgfplots, tcolorbox
% Beamer makes tens of thousands of \colorlet copies, and xcolor parses each one in full. Three kinds are shortcut, each
% giving exactly the definition xcolor would: a copy of a defined color by its plain name (a \let), a copy of the
% current color "." (PGF's xxcolor makes one at every \color) and a color mix (red!20!blue). The last two keep the
% color an expression produced together with the definitions it came from, and use it only while those are the same
% (\ifx) and xcolor's settings are as they were when the shortcut was installed. The kind of each color expression is
% kept by its text. Everything else (other expressions, types or models, tracing, xcolor's prologue file, a different
% xcolor \colorlet) runs xcolor's own code.
\makeatletter
\def\clearsnap@xc@end{\clearsnap@xc@end}
\def\clearsnap@xc@expect[#1]#2[#3]#4%
 {\begingroup
  \edef\@@cls{#1}\XC@edef\@@nam{#2}\XC@sdef\@@mod{#3}\XC@edef\@@clr{#4}%
  \XC@info\@@clr\@@tmp\XC@@tstfalse
  \ifnum\@@tmp=\@ne\ifx\@@mod\@empty
    \ifx\@@cls\@empty
      \XC@@tsttrue
    \else
      \edef\@@tmp
       {\expandafter\expandafter\expandafter\@secondoffive
         \csname\@backslashchar color@\@@clr\endcsname}%
      \ifx\@@cls\@@tmp\XC@@tsttrue\fi
    \fi
  \fi\fi
  \ifXC@@tst
    \XC@logdef\@@nam\XC@c@l@rlet\@@nam\@@clr
  \else
    \extractcolorspec\@@clr\@@clr
    \ifx\@@mod\@empty\else
      \expandafter\convertcolorspec\@@clr\@@mod\@@clr
      \edef\@@clr{{\@@mod}{\@@clr}}\fi
    \edef\@@tmp{\noexpand\XC@definecolor[\@@cls]{\@@nam}\@@clr}%
  \fi
  \expandafter\endgroup\@@tmp\xglobal@stop}
\def\clearsnap@xc@settings{%
  \ifconvertcolorsD D\fi\ifmaskcolors M\fi\ifblendcolors B\fi|%
  \expandafter\strip@prefix\meaning\XC@tgt@mod|\XC@tracing|\expandafter\strip@prefix\meaning\colorblend|%
  \ifx\rangeHsb\undefined\else\rangeHsb\fi|\ifx\rangetHsb\undefined\else\rangetHsb\fi|%
  \ifx\rangeRGB\undefined\else\rangeRGB\fi|\ifx\rangeHSB\undefined\else\rangeHSB\fi|%
  \ifx\rangeGray\undefined\else\rangeGray\fi|\ifx\XC@@xcp@\undefined\else\XC@@xcp@\fi}
% Installed after xcolor loads, if its \colorlet is the code above and it isn't tracing or writing a prologue.
\def\clearsnap@xc@install{%
  \ifx\clearsnap@xc@expect\XC@col@rlet
    \ifnum\XC@tracing<\tw@ \ifnum\XC@@xcp@=\z@\else
      \global\let\clearsnap@xc@orig\XC@col@rlet
      \gdef\XC@col@rlet[##1]##2[##3]##4{%
        \let\clearsnap@xc@next\clearsnap@xc@orig
        \ifcat$\detokenize{##1##3}$\clearsnap@xc@look{##4}\fi
        \clearsnap@xc@next[##1]{##2}[##3]{##4}}%
    \fi\fi
  \fi}
% The kind of the expression, kept by its text while that text can be read without xcolor's own expansion (no active
% punctuation, the usual \escapechar), decides what runs: one of the shortcuts below or xcolor's code. Kinds, mixes
% and current colors are kept for the first 2000 (\clearsnap@xc@count), so a document that makes colors without end
% (a color map) keeps the rest in xcolor's code.
\newcount\clearsnap@xc@count
\def\clearsnap@xc@simple{%
  \ifnum\escapechar=92 \ifnum\catcode`\!=13 1\else\ifnum\catcode`\:=13 1\else\ifnum\catcode`\-=13 1\else
  \ifnum\catcode`\+=13 1\else\ifnum\catcode`\;=13 1\else\ifnum\catcode`\"=13 1\else\ifnum\catcode`\>=13 1\else
  0\fi\fi\fi\fi\fi\fi\fi\else1\fi}
\def\clearsnap@xc@look#1{%
  \ifnum\clearsnap@xc@simple=\z@
    \let\clearsnap@xc@nm\clearsnap@xc@nmfast
    \edef\clearsnap@xc@clr{#1}%
    \edef\clearsnap@xc@clr{\detokenize\expandafter{\clearsnap@xc@clr}}%
    \ifcsname clearsnap@xc@k@\clearsnap@xc@clr\endcsname
      \expandafter\let\expandafter\clearsnap@xc@kind\csname clearsnap@xc@k@\clearsnap@xc@clr\endcsname
    \else
      \clearsnap@xc@classify
      \ifnum\clearsnap@xc@count<2000
        \global\advance\clearsnap@xc@count\@ne
        \global\expandafter\let\csname clearsnap@xc@k@\clearsnap@xc@clr\endcsname\clearsnap@xc@kind
      \fi
    \fi
  \else
    \let\clearsnap@xc@nm\clearsnap@xc@nmslow
    \XC@edef\clearsnap@xc@clr{#1}%
    \clearsnap@xc@classify
  \fi
  \clearsnap@xc@kind}
\def\clearsnap@xc@nmfast#1{\edef\clearsnap@xc@nam{#1}\edef\clearsnap@xc@nam{\detokenize\expandafter{\clearsnap@xc@nam}}}
\def\clearsnap@xc@nmslow#1{\XC@edef\clearsnap@xc@nam{#1}}
% A mix (has !), the current color, a plain name (xcolor's type 1: no >, :, ! or leading -, not .) or anything else.
\def\clearsnap@xc@classify{%
  \let\clearsnap@xc@kind\@empty
  \expandafter\in@\expandafter!\expandafter{\clearsnap@xc@clr}%
  \ifin@
    \let\clearsnap@xc@kind\clearsnap@xc@mix
  \else\ifx\clearsnap@xc@clr\clearsnap@xc@dot
    \let\clearsnap@xc@kind\clearsnap@xc@current
  \else
    \expandafter\in@\expandafter>\expandafter{\clearsnap@xc@clr}%
    \ifin@\else
      \expandafter\in@\expandafter:\expandafter{\clearsnap@xc@clr}%
      \ifin@\else
        \if-\expandafter\@car\clearsnap@xc@clr\relax\@nil\else
          \edef\clearsnap@xc@tmp{\noexpand\in@{/./}{/\clearsnap@xc@clr/}}\clearsnap@xc@tmp
          \ifin@\else\let\clearsnap@xc@kind\clearsnap@xc@plain\fi
        \fi
      \fi
    \fi
  \fi\fi}
\def\clearsnap@xc@mix{%
  \edef\clearsnap@xc@now{\clearsnap@xc@settings}%
  \ifx\clearsnap@xc@now\clearsnap@xc@set
    \ifcsname clearsnap@xc@r@\clearsnap@xc@clr\endcsname
      \let\clearsnap@xc@same\@empty
      \expandafter\expandafter\expandafter\clearsnap@xc@check
        \csname clearsnap@xc@n@\clearsnap@xc@clr\endcsname\clearsnap@xc@end
      \ifx\clearsnap@xc@same\@empty\let\clearsnap@xc@next\clearsnap@xc@hit\else\clearsnap@xc@changed\fi
    \else
      \ifnum\clearsnap@xc@count<2000 \let\clearsnap@xc@next\clearsnap@xc@keep\fi
    \fi
  \fi}
% A mix whose colors keep changing (pgfplots' "mapped color!80!black" changes at every patch of a surface) is kept
% again three times, then left to xcolor's code.
\def\clearsnap@xc@changed{%
  \ifcsname clearsnap@xc@v@\clearsnap@xc@clr\endcsname
    \expandafter\xdef\csname clearsnap@xc@v@\clearsnap@xc@clr\endcsname
      {\the\numexpr\csname clearsnap@xc@v@\clearsnap@xc@clr\endcsname+1}%
  \else
    \expandafter\gdef\csname clearsnap@xc@v@\clearsnap@xc@clr\endcsname{1}%
  \fi
  \ifnum\csname clearsnap@xc@v@\clearsnap@xc@clr\endcsname<4
    \let\clearsnap@xc@next\clearsnap@xc@keep
  \else
    \global\expandafter\let\csname clearsnap@xc@k@\clearsnap@xc@clr\endcsname\@empty
  \fi}
% A copy of a defined color by its plain name is the \let xcolor would make.
\def\clearsnap@xc@plain{%
  \ifnum\XC@tracing<\tw@
    \ifcsname\@backslashchar color@\clearsnap@xc@clr\endcsname
      \expandafter\ifx\csname\@backslashchar color@\clearsnap@xc@clr\endcsname\relax\else
        \let\clearsnap@xc@next\clearsnap@xc@copy
      \fi
    \fi
  \fi}
\def\clearsnap@xc@copy[#1]#2[#3]#4{%
  \clearsnap@xc@nm{#2}%
  \ifglobalcolors\global\else\xglobal@\fi
  \expandafter\let\csname\@backslashchar color@\clearsnap@xc@nam\expandafter\endcsname
    \csname\@backslashchar color@\clearsnap@xc@clr\endcsname
  \xglobal@stop}
% A copy of the current color "." (PGF's xxcolor makes one at every \color, so Beamer makes thousands) only depends on
% the current color's definition and xcolor's settings, unless that is a named-model color, so the copy each
% definition gave is kept by its meaning and used while the definition is the same (\ifx).
\def\clearsnap@xc@current{%
  \ifcsname\@backslashchar color@.\endcsname
    \ifnum\XC@type{.}>\@ne
      \edef\clearsnap@xc@now{\clearsnap@xc@settings}%
      \ifx\clearsnap@xc@now\clearsnap@xc@set
        \edef\clearsnap@xc@key{\expandafter\expandafter\expandafter\strip@prefix
          \expandafter\meaning\csname\@backslashchar color@.\endcsname}%
        \ifcsname clearsnap@xc@ps@\clearsnap@xc@key\endcsname
          \expandafter\ifx\csname clearsnap@xc@ps@\clearsnap@xc@key\expandafter\endcsname
              \csname\@backslashchar color@.\endcsname
            \let\clearsnap@xc@next\clearsnap@xc@curhit
          \else
            \clearsnap@xc@curnew
          \fi
        \else
          \ifnum\clearsnap@xc@count<2000 \clearsnap@xc@curnew\fi
        \fi
      \fi
    \fi
  \fi}
\def\clearsnap@xc@curnew{%
  \expandafter\expandafter\expandafter\clearsnap@xc@curmod\csname\@backslashchar color@.\endcsname\clearsnap@xc@end
  \ifx\clearsnap@xc@tmp\XC@mod@named\else\let\clearsnap@xc@next\clearsnap@xc@curkeep\fi}
\def\clearsnap@xc@curmod\xcolor@#1#2#3#4\clearsnap@xc@end{\XC@sdef\clearsnap@xc@tmp{#3}}
\def\clearsnap@xc@curhit[#1]#2[#3]#4{%
  \clearsnap@xc@nm{#2}%
  \ifglobalcolors\global\else\xglobal@\fi
  \expandafter\let\csname\@backslashchar color@\clearsnap@xc@nam\expandafter\endcsname
    \csname clearsnap@xc@pr@\clearsnap@xc@key\endcsname
  \xglobal@stop}
\def\clearsnap@xc@curkeep[#1]#2[#3]#4{%
  \global\advance\clearsnap@xc@count\@ne
  \let\clearsnap@xc@ckey\clearsnap@xc@key
  \expandafter\let\expandafter\clearsnap@xc@src\csname\@backslashchar color@.\endcsname
  \clearsnap@xc@orig[#1]{#2}[#3]{#4}%
  \clearsnap@xc@nm{#2}%
  \global\expandafter\let\csname clearsnap@xc@pr@\clearsnap@xc@ckey\expandafter\endcsname
    \csname\@backslashchar color@\clearsnap@xc@nam\endcsname
  \global\expandafter\let\csname clearsnap@xc@ps@\clearsnap@xc@ckey\endcsname\clearsnap@xc@src}
\def\clearsnap@xc@check#1{%
  \def\clearsnap@xc@tmp{#1}%
  \ifx\clearsnap@xc@tmp\clearsnap@xc@endname\else
    \expandafter\ifx\csname\@backslashchar color@#1\expandafter\endcsname
      \csname clearsnap@xc@d@\clearsnap@xc@clr @#1\endcsname
    \else\def\clearsnap@xc@same{no}\fi
    \expandafter\clearsnap@xc@check
  \fi}
\def\clearsnap@xc@endname{\clearsnap@xc@end}
\def\clearsnap@xc@hit[#1]#2[#3]#4{%
  \clearsnap@xc@nm{#2}%
  \ifglobalcolors\global\else\xglobal@\fi
  \expandafter\let\csname\@backslashchar color@\clearsnap@xc@nam\expandafter\endcsname
    \csname clearsnap@xc@r@\clearsnap@xc@clr\endcsname
  \xglobal@stop}
% Finds the colors a mix names (name!num!name...[!num], one leading - allowed), keeps their definitions, lets xcolor
% compute the color and keeps it. A mix that names anything else (the current color ".", extended syntax, an
% undefined color) is only computed.
\def\clearsnap@xc@keep[#1]#2[#3]#4{%
  \let\clearsnap@xc@names\@empty \let\clearsnap@xc@good\relax
  \expandafter\clearsnap@xc@name\clearsnap@xc@clr!\clearsnap@xc@end
  \ifx\clearsnap@xc@good\relax
    \global\advance\clearsnap@xc@count\@ne
    \let\clearsnap@xc@key\clearsnap@xc@clr
    \def\clearsnap@xc@one##1{\expandafter\let\csname clearsnap@xc@t@##1\expandafter\endcsname\csname\@backslashchar color@##1\endcsname}%
    \expandafter\clearsnap@xc@each\clearsnap@xc@names\clearsnap@xc@end
    \clearsnap@xc@orig[#1]{#2}[#3]{#4}%
    \clearsnap@xc@nm{#2}%
    \def\clearsnap@xc@one##1{\global\expandafter\let\csname clearsnap@xc@d@\clearsnap@xc@key @##1\expandafter\endcsname\csname clearsnap@xc@t@##1\endcsname}%
    \expandafter\clearsnap@xc@each\clearsnap@xc@names\clearsnap@xc@end
    \global\expandafter\let\csname clearsnap@xc@n@\clearsnap@xc@key\endcsname\clearsnap@xc@names
    \global\expandafter\let\csname clearsnap@xc@r@\clearsnap@xc@key\expandafter\endcsname
      \csname\@backslashchar color@\clearsnap@xc@nam\endcsname
  \else
    \clearsnap@xc@orig[#1]{#2}[#3]{#4}%
  \fi}
\def\clearsnap@xc@each#1{%
  \def\clearsnap@xc@tmp{#1}%
  \ifx\clearsnap@xc@tmp\clearsnap@xc@endname\else\clearsnap@xc@one{#1}\expandafter\clearsnap@xc@each\fi}
\def\clearsnap@xc@name#1!#2\clearsnap@xc@end{%
  \clearsnap@xc@isname{#1}%
  \def\clearsnap@xc@tmp{#2}%
  \ifx\clearsnap@xc@tmp\@empty\else\clearsnap@xc@num#2\clearsnap@xc@end\fi}
\def\clearsnap@xc@num#1!#2\clearsnap@xc@end{%
  \clearsnap@xc@isnum{#1}%
  \def\clearsnap@xc@tmp{#2}%
  \ifx\clearsnap@xc@tmp\@empty\else\clearsnap@xc@name#2\clearsnap@xc@end\fi}
\def\clearsnap@xc@isnum#1{%
  \def\clearsnap@xc@tmp{#1}%
  \ifx\clearsnap@xc@tmp\@empty\let\clearsnap@xc@good\@empty\else
    \in@{ }{#1}\ifin@\let\clearsnap@xc@good\@empty\else\clearsnap@xc@digit#1\clearsnap@xc@end\fi
  \fi}
\def\clearsnap@xc@digit#1{%
  \ifx#1\clearsnap@xc@end\else
    \ifnum`#1<`0 \ifnum`#1=`.\else\let\clearsnap@xc@good\@empty\fi\else
      \ifnum`#1>`9 \let\clearsnap@xc@good\@empty\fi\fi
    \expandafter\clearsnap@xc@digit
  \fi}
\def\clearsnap@xc@isname#1{%
  \def\clearsnap@xc@tmp{#1}%
  \ifx\clearsnap@xc@tmp\@empty\let\clearsnap@xc@good\@empty\else
    \clearsnap@xc@unplain{#1}%
    \ifx\clearsnap@xc@good\relax
      \edef\clearsnap@xc@tmp{\clearsnap@xc@nodash#1\clearsnap@xc@end}%
      \let\clearsnap@xc@def\@empty
      \ifcsname\@backslashchar color@\clearsnap@xc@tmp\endcsname
        \expandafter\ifx\csname\@backslashchar color@\clearsnap@xc@tmp\endcsname\relax\else
          \let\clearsnap@xc@def\relax
        \fi
      \fi
      \ifx\clearsnap@xc@def\relax
        \edef\clearsnap@xc@names{\clearsnap@xc@names{\clearsnap@xc@tmp}}%
      \else
        \let\clearsnap@xc@good\@empty
      \fi
    \fi
  \fi}
\def\clearsnap@xc@nodash#1#2\clearsnap@xc@end{\if-#1#2\else#1#2\fi}
\def\clearsnap@xc@dot{.}
\def\clearsnap@xc@unplain#1{%
  \in@{:}{#1}\ifin@\let\clearsnap@xc@good\@empty\fi
  \in@{>}{#1}\ifin@\let\clearsnap@xc@good\@empty\fi
  \in@{,}{#1}\ifin@\let\clearsnap@xc@good\@empty\fi
  \in@{;}{#1}\ifin@\let\clearsnap@xc@good\@empty\fi
  \in@{/}{#1}\ifin@\let\clearsnap@xc@good\@empty\fi
  \in@{=}{#1}\ifin@\let\clearsnap@xc@good\@empty\fi
  \in@{+}{#1}\ifin@\let\clearsnap@xc@good\@empty\fi
  \in@{--}{#1}\ifin@\let\clearsnap@xc@good\@empty\fi
  \def\clearsnap@xc@tmp{#1}%
  \ifx\clearsnap@xc@tmp\clearsnap@xc@dot\let\clearsnap@xc@good\@empty\fi
  \edef\clearsnap@xc@tmp{\clearsnap@xc@nodash#1\clearsnap@xc@end}%
  \ifx\clearsnap@xc@tmp\clearsnap@xc@dot\let\clearsnap@xc@good\@empty\fi}
\def\clearsnap@xc@start{\clearsnap@xc@install\xdef\clearsnap@xc@set{\clearsnap@xc@settings}}
\@ifpackageloaded{xcolor}{\clearsnap@xc@start}{\AddToHook{package/xcolor/after}{\clearsnap@xc@start}}
\makeatother
